\uuid{BUzn}
\titre{Équations homogènes/solution générale}
\auteur{hanani ; mzari}
\organisation{exo7}
\licence{CC BY-NC-SA 4.0}
\datecreate{2020-01-20}
\niveau{L1}
\module{Analyse}
\chapitre{Equation différentielle}
\sousChapitre{Résolution d'équation différentielle du premier ordre}
\difficulte{1}
\theme{Equations du premier ordre, équations homogènes, solution générale}
\exo7qcmid{lille-2.equations-differentielles.002}

\contenu{
\begin{question}
\qclassification{225.01}
\qtags{motcle=équations homogènes/solution générale}
On considère les équations différentielles :
$$(E_1)\; :\; (1+x^2)y'-y=0\quad\mbox{et}\quad (E_2)\; :\; y'-\frac{y}{1+x^2}=0.$$
Parmi les affirmations suivantes cocher celles qui sont vraies :
\begin{answers}
\bad{Si $y$ est une solution de $(E_1)$, alors $\displaystyle z=\frac{y}{1+x^2}$ est une solution de $(E_2)$.}
\good{$(E_1)$ est $(E_2)$ ont les mêmes solutions.}
\good{La solution générale de $(E_1)$ sur $\Rr$ est : $\displaystyle y=k\mathrm{e}^{arctan (x)}$, $k\in \Rr$.}
\bad{La solution générale de $(E_2)$ sur $\Rr$ est : $\displaystyle y=k\arctan (x)$, $k\in \Rr$.}
\end{answers}
\begin{explanations}
En divisant $(E_1)$ par $1+x^2$, on obtient $(E_2)$ : $(E_1)$ et $(E_2)$ ont les mêmes solutions. Une primitive de $\displaystyle \frac{-1}{1+x^2}$ est $-\arctan x$, donc la solution générale de $(E_1)$ sur $\Rr$ est : $y=k\mathrm{e}^{arctan (x)}$, $k\in \Rr$.
\end{explanations}
\end{question}
}
